% Chapter 2 entries ATTRIB .. DISKCOPY (AN-01).  Sources per entry in the comment
% above it; every Example is a capture listed in screens/SCREENS.

% ATTRIB -- os/milton/dosutil.c attrib_main/attrib_walk (+R -R +A -A /S,
% display R col 1 / A col 4 / full path col 10; messages); shipped as
% APPS\ATTRIB.COM.  Gates: test-attrib, test-dosutil-emu.
\begin{DosEntry}{ATTRIB}{Attribute}{external}
\Purpose Sets or displays the read-only and archive attributes of a file.
\Format \SyntaxLine{\opt{\ph{d}\kw{:}}\opt{\ph{path}}\kw{ATTRIB}
  \opt{\kw{+R}\texttt{|}\kw{-R}} \opt{\kw{+A}\texttt{|}\kw{-A}}
  \opt{\ph{d}\kw{:}}\opt{\ph{path}}\ph{filename}\opt{\kw{.}\ph{ext}}
  \opt{\kw{/S}}}
\Type{External}
\Remarks
\Param{+R} makes the file read-only.  A read-only file can be read, typed
and copied, but it cannot be changed or deleted.

\Param{-R} makes the file read-write again.

\Param{+A} sets the archive attribute of the file.  The archive attribute
marks a file that has been created or changed since it was last backed up;
\cmd{XCOPY} uses it with its \kw{/A} and \kw{/M} parameters.

\Param{-A} clears the archive attribute.

\Param{filename} names the files to change or display.  Wildcards are
allowed.  Directories are not changed.

\Param{/S} processes the matching files in the named directory and in all
of its subdirectories.

If you give no \kw{+} or \kw{-} parameter, \cmd{ATTRIB} displays the
attributes of each matching file: \texttt{R} in the first column for a
read-only file, \texttt{A} in the fourth column when the archive attribute is
set, then the full path of the file.

If no file matches, or the name is a directory, \cmd{ATTRIB} displays
\texttt{Invalid path (or file not found)}.  If you give no file name, or
more than one, it displays \texttt{Invalid number of parameters}.

\cmd{ATTRIB} returns exit code 0 when it succeeds and 1 when it displays a
message.
\Example
The first command makes \file{MEMO.TXT} read-only; the second displays the
attributes of every file with the extension \file{TXT}; the \cmd{DEL}
command then cannot delete the read-only file:
\Shot{attrib}
To display the attributes of the program files in every directory of the
disk, begin at the root directory and use \kw{/S}:
\Shot{attribs}
\end{DosEntry}

% BREAK -- os/milton/command.c builtin_break (ON/OFF case-insensitive, other
% arguments leave the state; "BREAK is on|off"); sysinit.c (default ON
% unless CONFIG.SYS BREAK=OFF); int21.c conin_cooked_line (CTRL-C on CON).
% ADR-0003 Amendment DEC-16.  Gates: test-command, test-j1lu-control-emu.
\begin{DosEntry}{BREAK}{Control-Break}{internal}
\Purpose Sets or displays the CTRL-BREAK check.
\Format \SyntaxLine{\kw{BREAK} \opt{\kw{ON}\texttt{|}\kw{OFF}}}
\Type{Internal}
\Remarks
Type \cmd{BREAK} with no parameter to display the current setting:
\texttt{BREAK is on} or \texttt{BREAK is off}.  After \kw{ON} or \kw{OFF},
\cmd{BREAK} also displays the new setting.

The check is on when InitechDOS starts, unless \file{CONFIG.SYS} contains
the configuration command \texttt{BREAK=OFF} (Chapter~3).  Any parameter
other than \kw{ON} or \kw{OFF} leaves the setting unchanged.

Whatever the setting, InitechDOS checks for \key{Ctrl-C} whenever it reads
a line from the keyboard: the line is discarded, a program reading it is
stopped, and the command prompt is displayed again.
\draft{see bead initech-vlhp (DEC-16 C-6): in DOS 3.3, BREAK ON also checks for CTRL-C during disk
and other operations; the shipped build records and reports the setting
but checks only during keyboard input}
With \kw{ON}, InitechDOS also checks for \key{Ctrl-C} while a program reads
or writes the disk, so a long operation can be stopped.
\Example
\Shot{break}
\end{DosEntry}

% CD -- os/milton/command.c builtin_cd (no arg: cwd_display; failure:
% MSG-DOS-0018 "Invalid directory"); int21.c do_chdir (relative, absolute,
% "..", "\" alone = the root since initech-33t2).  Gates: test-ut6d,
% test-fat-subdir, test-shell.
\begin{DosEntry}{CD}{Change Directory}{internal}
\AlsoNamed{CHDIR}
\Purpose Changes the current directory, or displays its name.
\Format \SyntaxLine{\kw{CD} \opt{\ph{d}\kw{:}}\opt{\ph{path}}}
  \SyntaxLine{\kw{CHDIR} \opt{\ph{d}\kw{:}}\opt{\ph{path}}}
\Type{Internal}
\Remarks
\Param{path} names the directory that is to become the current directory.
A path that begins with a backslash starts at the root directory; any other
path starts at the current directory.  \file{..} names the parent of the
current directory.

Type \cmd{CD} with no parameter to display the drive and full path of the
current directory.

The command prompt shows the current directory, so a successful \cmd{CD}
displays nothing more.  If the directory does not exist, \cmd{CD} displays
\texttt{Invalid directory} and the current directory does not change.

\type{CD \bs} makes the root directory the current directory, from any
directory on the disk.
\Example
The example moves into \file{SAMPLES}, displays the current directory,
moves back with \type{CD ..}, moves in again and goes straight to the root
with \type{CD \bs}:
\Shot{cd}
\end{DosEntry}

% CHKDSK -- os/milton/chkdsk.c (parse: one file operand, /F /V; report
% lines; contiguity; lost chains + FILE0000.CHK; cross-links reported, never
% repaired; memory from the live MCB chain); chkdsk_program.c.  Shipped as
% DOS\CHKDSK.COM.  Gates: test-chkdsk, test-chkdsk-emu.
\begin{DosEntry}{CHKDSK}{Check Disk}{external}
\Purpose Checks the directories and the file allocation table of a disk and
reports on the disk and on memory.
\Format \SyntaxLine{\opt{\ph{d}\kw{:}}\opt{\ph{path}}\kw{CHKDSK}
  \opt{\ph{d}\kw{:}}\opt{\ph{path}}\opt{\ph{filename}\opt{\kw{.}\ph{ext}}}
  \opt{\kw{/F}} \opt{\kw{/V}}}
\Type{External}
\Remarks
\cmd{CHKDSK} reads every directory of the disk and follows the chain of
allocation units of every file.  It then displays a status report:
\begin{itemize}
\item the volume label and the date and time the volume was created
\item the total disk space, in bytes
\item the bytes and number of hidden files, of directories, and of user
  files
\item the bytes still available on the disk
\item the total memory and the free memory, in bytes
\end{itemize}

\Param{filename} names files to check for fragmentation.  Wildcards are
allowed.  For each named file stored in more than one piece, \cmd{CHKDSK}
displays its path and \texttt{Contains} \ph{n} \texttt{non-contiguous
blocks}; when every named file is in one piece it displays \texttt{All
specified file(s) are contiguous}.  If no file matches, it displays
\texttt{File not found}.

\Param{/F} lets \cmd{CHKDSK} correct the errors it finds.  Without
\kw{/F}, \cmd{CHKDSK} reports what it would do and displays
\texttt{Corrections will not be written to disk}; nothing on the disk is
changed.

\Param{/V} displays the full path of every file as it is checked.

When \cmd{CHKDSK} finds allocation units that are marked in use but belong
to no file, it displays
\ph{n} \texttt{lost allocation units found in} \ph{n} \texttt{chains.}
and asks \texttt{Convert lost chains to files (Y/N)?}  With \kw{/F}, answer
\type{Y} to save each chain as a file named \file{FILE0000.CHK},
\file{FILE0001.CHK}, and so on, in the root directory, where you can examine
it; answer \type{N} to free the space.  Without \kw{/F}, \cmd{CHKDSK}
displays the number of bytes that would be freed.

A file whose allocation chain points outside the disk is reported with
\texttt{Has invalid cluster, file truncated}.  Two files that share an
allocation unit are both reported with \texttt{File is cross-linked: on
cluster} \ph{n}; \cmd{CHKDSK} does not repair a cross-link.  Copy both
files, check their contents, and delete the damaged one.

\cmd{CHKDSK} returns exit code 0 when the disk has no errors and 1 when it
reports errors.
\Example
To check the practice diskette:
\Shot{chkdsk}
To check whether \file{MEMO.TXT} is stored in one piece:
\Shot{chkfile}
\end{DosEntry}

% CLS -- os/milton/command.c builtin_cls (form feed: clear and home; any
% parameters ignored).  Gate: test-command, test-pubs-screens (every screen
% begins with CLS).
\begin{DosEntry}{CLS}{Clear Screen}{internal}
\Purpose Clears the display.
\Format \SyntaxLine{\kw{CLS}}
\Type{Internal}
\Remarks
The display is cleared and the command prompt appears on the top line.
Parameters typed after \cmd{CLS} are ignored.
\Example
Type \type{CLS} and press \key{Enter}.  The screen is cleared and shows only
the command prompt, \texttt{A:\bs>}.  Every example in this book begins
this way.
\end{DosEntry}

% COPY -- os/milton/command.c builtin_copy (two operands; directory target
% keeps the name; same-file refusal MSG-DOS-0020 + "0 File(s) copied";
% source missing MSG-DOS-0003; create refused MSG-DOS-0009 or MSG-DOS-0002;
% footer "        1 File(s) copied", MSG-DOS-0061 since WL-0154); COPY CON
% ends at Ctrl-Z (WL-0102).
% Gates: test-copy-selfcopy, test-dos-safety-k01, test-dos-safety-k02,
% test-j1lu-control-emu.  Date stamp of the copy: initech-dwbl.
\begin{DosEntry}{COPY}{Copy File}{internal}
\Purpose Copies a file.
\Format \SyntaxLine{\kw{COPY} \opt{\ph{d}\kw{:}}\opt{\ph{path}}\ph{filename}\opt{\kw{.}\ph{ext}}
  \opt{\ph{d}\kw{:}}\opt{\ph{path}}\opt{\ph{filename}\opt{\kw{.}\ph{ext}}}}
\Type{Internal}
\Remarks
The first file specification names the \emph{source}, the file to copy.
The second names the \emph{target}: the new file, or the directory to copy
into.  Both are required.

If the target is a directory, the copy is put in that directory with the
name of the source file.  If the target is an existing file, it is
replaced.  When the copy is complete, \cmd{COPY} displays
\texttt{1 File(s) copied}.\draft{see bead initech-bc2j: the shipped COPY
also accepts wildcards, combines files named with +, takes the /A, /B and
/V switches, and copies into the current directory when no target is given;
this entry is to be rewritten}

To copy a directory and its subdirectories, use \cmd{XCOPY}.

You can copy from the keyboard by naming the device \file{CON} as the
source: type the lines of the file, then press \key{Ctrl-Z} and
\key{Enter}.  \texttt{COPY CON NOTE.TXT} creates the file \file{NOTE.TXT}
from what you type.

If the source and target are the same file, \cmd{COPY} displays
\texttt{File cannot be copied onto itself} and \texttt{0 File(s) copied},
and the file is not changed.  If the source does not exist, it displays
\texttt{File not found}.  If the target is read-only, it displays
\texttt{Access denied}; if the target cannot be created for another reason,
such as a full disk, it displays \texttt{Bad command or file name}.

\draft{bead needed (follow-up of initech-dwbl): the shipped COPY stamps the
copy with the current date and time; DOS 3.3 keeps the source's date and
time}
The copy has the same date and time as the source file.
\Example
The first command copies \file{MEMO.TXT} to a new file, \file{MEMO.BAK}; the
second copies \file{PHONE.TXT} into the directory \file{SAMPLES}; the third
names the same file twice:
\Shot{copy}
\end{DosEntry}

% DATE -- os/milton/command.c builtin_date, cmd_parse_date (mm-dd-yy[yy],
% '-' or '/' used twice; 80..99 -> 19xx, 00..79 -> 20xx; four-digit years from 1980 on, upper bound in cmd_parse_date;
% days per month with leap years), cmd_format_date ("Current date is Day
% mm-dd-yyyy"); "Invalid date" then the unchanged date.  No prompt for a new
% date in the shipped build (documented as built).  Gates: test-datetime,
% test-command.
\begin{DosEntry}{DATE}{}{internal}
\Purpose Sets or displays the system date.
\Format \SyntaxLine{\kw{DATE} \opt{\ph{mm}\kw{-}\ph{dd}\kw{-}\ph{yy}}}
\Type{Internal}
\Remarks
\Param{mm} is the month, 1 through 12.

\Param{dd} is the day, 1 through the last day of the month.

\Param{yy} is the year.  Two digits from 80 through 99 stand for the years
from 1980 to the end of the century; two digits from 00 through 79 stand for
the years of the next century.  You may also give all four digits.

Separate the parts with hyphens or with slashes, the same character both
times.  InitechDOS works out the day of the week for you.

With a valid date, \cmd{DATE} sets the system date and displays it.  With
no parameter, it displays the current date:
\texttt{Current date is} followed by the day of the week and the date.

If the date is not valid (for example, February~30), \cmd{DATE} displays
\texttt{Invalid date}, leaves the system date unchanged and displays it.

\Example
\Shot{date}
\end{DosEntry}

% DEL -- os/milton/command.c builtin_del (wildcards: restart FINDFIRST per
% unlink; *.* asks MSG-DOS-0012 first, only Y deletes; read-only:
% MSG-DOS-0009 + the names left; MSG-DOS-0003; MSG-DOS-0011).  Gates:
% test-dos-safety-k03, test-dos-safety-k04, test-dos-safety-k15,
% test-dos-safety-del-failure.
\begin{DosEntry}{DEL}{Delete}{internal}
\AlsoNamed{ERASE}
\Purpose Deletes files.
\Format \SyntaxLine{\kw{DEL} \opt{\ph{d}\kw{:}}\opt{\ph{path}}\ph{filename}\opt{\kw{.}\ph{ext}}}
  \SyntaxLine{\kw{ERASE} \opt{\ph{d}\kw{:}}\opt{\ph{path}}\ph{filename}\opt{\kw{.}\ph{ext}}}
\Type{Internal}
\Remarks
\Param{filename} names the file or files to delete.  Wildcards are
allowed: \type{DEL *.BAK} deletes every file with the extension \file{BAK}
in the current directory.

To delete every file in a directory, name \file{*.*}.  \cmd{DEL} first asks
\texttt{Are you sure (Y/N)?}  Type \type{Y} and press \key{Enter} to delete
the files; any other answer deletes nothing.

\cmd{DEL} displays nothing when it succeeds.  If no file matches, it
displays \texttt{File not found}.  A read-only file is not deleted:
\cmd{DEL} displays \texttt{Access denied}, and when you used wildcards it
also lists the matching files that remain.  \cmd{DEL} does not delete directories; use \cmd{RD}.

\begin{Important}
A deleted file cannot be recovered with InitechDOS commands.
\end{Important}
\Example
The example makes a copy of \file{MEMO.TXT}, refuses to delete every file,
deletes the copy, and shows that it is gone:
\Shot{del}
\end{DosEntry}

% DIR -- os/milton/command.c builtin_dir (MSG-DOS-0021..0026 display
% fragments, cmd_format_dir_line, footer from AH=36h; empty -> File not
% found; e47v: one operand -- drive A only, else MSG-DOS-0004 -- a directory
% or a file name, "name" = name.* and ".ext" = *.ext, /P pauses with
% MSG-DOS-0016 after 23 rows, /W five names to a line, a subdirectory in
% brackets; any other switch or a second operand MSG-DOS-0010).
% Gates: test-dir, test-dir-look, test-dir-volume, test-dir-free-space,
% test-dir-page, test-dir-args-mutant (initech-e47v).
\begin{DosEntry}{DIR}{Directory}{internal}
\Purpose Lists the files and subdirectories in a directory.
\Format \SyntaxLine{\kw{DIR} \opt{\ph{d}\kw{:}}\opt{\ph{path}}\opt{\ph{filename}\opt{\kw{.}\ph{ext}}}
  \opt{\kw{/P}} \opt{\kw{/W}}}
\Type{Internal}
\Remarks
\Param{d:} is the drive; it can only be drive~A.  Any other drive displays
\texttt{Invalid drive specification}.

\Param{path} names the directory to list.  If you leave it out, \cmd{DIR}
lists the current directory.  Listing a directory does not change the
current directory.

\Param{filename} limits the listing to matching files.  Wildcards are
allowed.  A name without an extension lists every extension of that name
(\type{DIR MEMO} lists \file{MEMO.*}), and an extension alone lists every
name with that extension (\type{DIR .TXT} lists \file{*.TXT}).

\Param{/P} pauses when the screen is full and displays
\texttt{Press any key to continue . . .}; press any key to see the next
screen.

\Param{/W} lists the names only, five to a line, with the name of a
subdirectory in brackets.

Only one file specification may be given; a second one, or any switch
other than \kw{/P} and \kw{/W}, displays \texttt{Too many parameters}.

The listing begins with the volume label of the disk and the name of the
directory.  Each file then has one line: its name, its extension, its size
in bytes, and the date and time it was last changed, with \texttt{a} for
morning and \texttt{p} for afternoon.  A subdirectory shows
\texttt{<DIR>} in place of a size.  The last line gives the number of
entries listed and the number of bytes free on the disk.

If the directory has no entries, \cmd{DIR} displays the headings followed
by \texttt{File not found}.
\Example
To list the root directory of the practice diskette:
\Shot{dir}
To list only the text files, and then the directory \file{SAMPLES} in the
wide form:
\Shot{dirargs}
\end{DosEntry}

% DISKCOMP / DISKCOPY -- os/milton/diskutil.c copy_run (the single-drive
% swap flow of MS-DOS 3.3 UG p.26 and pp.59-63: FIRST/SECOND and
% SOURCE/TARGET prompts, 64 sectors per pass; "Compare error on disk / side
% s, track t", "Compare OK"; removable drives only, two-line refusal p.358;
% B: is MSG-DOS-0004 by DEC-18, spec/drive_letters.json; /1 /8 options;
% "Copy/Compare another diskette (Y/N)?").  WL-0147.  Shipped as
% DOS\DISKCOMP.COM and DOS\DISKCOPY.COM.  Gates: test-formatsys,
% test-formatsys-emu, test-diskcopy-emu.
\begin{DosEntry}{DISKCOMP}{Compare Diskette}{external}
\Purpose Compares the contents of two diskettes sector by sector.
\Format \SyntaxLine{\opt{\ph{d}\kw{:}}\opt{\ph{path}}\kw{DISKCOMP}
  \opt{\ph{d}\kw{:} \opt{\ph{d}\kw{:}}} \opt{\kw{/1}} \opt{\kw{/8}}}
\Type{External}
\Remarks
\Param{d:} names the drives of the first and second diskette.  If you give
none, or only one, drive~A is used for the others.  Both drives must be
diskette drives: for a fixed disk \cmd{DISKCOMP} displays
\texttt{Specified drive does not exist,} \texttt{or is non-removable}.  Drive~B
does not exist in this version (\texttt{Invalid drive specification}).

With one drive, \cmd{DISKCOMP} asks for the diskettes in turn:
\texttt{Insert FIRST diskette into drive A:} and \texttt{Insert SECOND
diskette into drive A:}.  Change the diskette and press any key.  It reads
up to 64 sectors of the first diskette at a time, so a large diskette is
compared in several passes.

\Param{/1} compares only the first side of each diskette.  \Param{/8}
compares only the first eight sectors of each track.

For every sector that differs, \cmd{DISKCOMP} displays
\texttt{Compare error on disk} and, on the next line, the side and track.
If no sector differs, it displays \texttt{Compare OK}.  It then asks
\texttt{Compare another diskette (Y/N)?}; type \type{Y} or \type{N}.
\Example
To compare two diskettes in drive~A, type \type{DISKCOMP A: A:}, then
insert the first and second diskettes as \cmd{DISKCOMP} asks for them.
\end{DosEntry}

\begin{DosEntry}{DISKCOPY}{Copy Diskette}{external}
\Purpose Copies the contents of one diskette onto another, sector by sector.
\Format \SyntaxLine{\opt{\ph{d}\kw{:}}\opt{\ph{path}}\kw{DISKCOPY}
  \opt{\ph{d}\kw{:} \opt{\ph{d}\kw{:}}} \opt{\kw{/1}}}
\Type{External}
\Remarks
\Param{d:} names the drives of the source and target diskette.  If you give
none, or only one, drive~A is used for the others.  Both drives must be
diskette drives: for a fixed disk \cmd{DISKCOPY} displays
\texttt{Specified drive does not exist,} \texttt{or is non-removable}.  Drive~B
does not exist in this version (\texttt{Invalid drive specification}).

With one drive, \cmd{DISKCOPY} asks for the diskettes in turn:
\texttt{Insert SOURCE diskette in drive A:} and \texttt{Insert TARGET
diskette into drive A:}, each followed by \texttt{Press any key when ready
. . .}  It reads up to 64 sectors of the source at a time, so a large
diskette is copied in several passes.  Everything on the target diskette is
replaced.

\Param{/1} copies only the first side of the diskette.

If the source cannot be read, \cmd{DISKCOPY} displays \texttt{FIRST
diskette bad or incompatible}; if the target cannot be written, it displays
\texttt{Target diskette may be unusable}.  When the copy is complete it asks
\texttt{Copy another diskette (Y/N)?}; type \type{Y} or \type{N}.
\Example
To copy a diskette using drive~A only, type \type{DISKCOPY A: A:}, then
insert the source and target diskettes as \cmd{DISKCOPY} asks for them.
\end{DosEntry}
